\chapter{Conclusion}
\label{chap:conclusion}

This thesis studied a rolling-horizon warehouse-distribution planning problem in which customer orders become visible over time, many orders can be served within flexible multi-day windows, and same-day execution is constrained not only by routing feasibility but also by warehouse-side resource limits. The central claim of the thesis is that such systems should not be treated as static route-planning problems. Instead, they require coordinated admission control, route construction, and depot-aware smoothing under a multi-day operational horizon.

Methodologically, the thesis developed along two complementary lines. The first was a retained controller-centered rolling-horizon line, which established a disciplined paper-facing baseline for cross-day admission, protection, and risk-aware clipping. This line remains the core reference backbone of the thesis because it provides a stable mechanism progression and a defensible experimental story under pressure. The second line emerged from later diagnostics, which showed that some remaining failures are more naturally explained by localized depot-wave bottlenecks than by broad admission pressure alone. This led to an integrated controller-routing-repair architecture in which dynamic shadow-price signals are used to internalize bucket-level depot pressure more explicitly.

The empirical findings support this two-layer interpretation. The retained controller-centered line demonstrates that better rolling admission logic materially improves service outcomes under pressure, with the \texttt{v2} $\rightarrow$ \texttt{v4} $\rightarrow$ \texttt{v5} $\rightarrow$ \texttt{v6f} $\rightarrow$ \texttt{v6g} progression providing the central paper-facing evidence chain. The integrated solver line then adds a different type of insight. In West, the hardest large instance, the dominant problem is not all-day capacity shortage but a sharply concentrated morning picking wave around 05:30--06:15. The V8.2 dynamic-shadow-price variant achieves the best service--smoothing trade-off found in this line, reaching approximately 99.41\% service while reducing both total and morning picking overload relative to the earlier OR-V7 reference. On Herlev, the same mechanism also improves service, while on East it becomes clear that this type of smoothing is not universally dominant and must be interpreted in a regime-dependent way.

Taken together, these results suggest that the right architecture for rolling warehouse-distribution planning is neither a pure route-only optimizer nor a controller that ignores the temporal structure of depot operations. Instead, the most useful perspective is an integrated one in which controller decisions shape daily demand exposure, routing decisions determine how that exposure is realized operationally, and depot-aware signals help prevent localized execution waves from undermining service quality. Future work should therefore focus on stronger regime-aware activation, richer stochastic robustness, and tighter integration between routing decisions and labor-sensitive warehouse planning.
